\documentclass[11pt, a4paper]{article}

\usepackage[T1]{fontenc}
\usepackage[utf8]{inputenc}
\usepackage{tgpagella}
\usepackage[spanish, es-noshorthands]{babel}
\usepackage{amsmath, amssymb}
\usepackage[most]{tcolorbox}
\usepackage[american]{circuitikz}
\usepackage[margin=2.5cm]{geometry}
\usepackage{fancyhdr}

\pagestyle{fancy}
\fancyhf{}
\setlength{\headheight}{16pt}
\lhead{\textbf{UCB Sede Tarija} | Ing. Mecatrónica}
\rhead{IMT-221 | Ejercicios S06-3}
\renewcommand{\headrulewidth}{0.4pt}
\definecolor{ucbblue}{RGB}{0, 51, 102}

\begin{document}

\begin{tcolorbox}[colback=ucbblue!5, colframe=ucbblue, title=\textbf{Ejercicio Integrado: Fuente - Amplificador - Carga}]
    Aplica el modelo completo para determinar la ganancia global del sistema $G_v$.
\end{tcolorbox}

\section*{Circuito Equivalente AC}
\textbf{Datos asignados:} $I_{CQ}=1.0\,\text{mA}$, $\beta=100$, $R_B=100\,\text{k}\Omega$, $R_S=10\,\text{k}\Omega$, $R_C=4.7\,\text{k}\Omega$, $R_L=10\,\text{k}\Omega$, $T \approx 300\,\text{K}$ ($V_T=25.85\,\text{mV}$). Suponga $r_o \to \infty$.

\begin{center}
\begin{circuitikz}[scale=1, transform shape, thick]
    % Source part
    \draw (-2,0) to[sV, l={$v_s$}] (-2,2) to[R, l={$R_S$}] (0,2);
    \draw (-2,0) -- (0,0) node[ground]{};
    % Ampli Input
    \draw (0,2) to[short, *-] (2,2) node[right]{B};
    \draw (0,2) to[R, l_={$R_B$}] (0,0) node[ground]{};
    \draw (2,2) to[R, l={$r_\pi$}] (2,0) node[ground]{};
    % Ampli Output
    \draw (5,2) node[left]{C} to[cI, l_={$g_m v_\pi$}] (5,0) node[ground]{};
    \draw (5,2) -- (7,2);
    \draw (6,2) to[R, l={$R_C$}] (6,0) node[ground]{};
    \draw (7,2) to[R, l={$R_L$}] (7,0) node[ground]{};
    \draw (7,2) to[short, *-o] (8,2) node[right]{$v_o$};
\end{circuitikz}
\end{center}

\section*{Secuencia de Análisis}
\begin{enumerate}
    \item Calcula los parámetros intrínsecos $g_m$ y $r_\pi$.
    \item Calcula la resistencia de entrada total $R_{in}$.
    \item \textbf{Predicción:} ¿Será la tensión de entrada $v_i$ igual, menor o mayor que $v_s$? Calcula la atenuación de entrada $v_i / v_s$.
    \item \textbf{Predicción:} ¿Qué efecto tendrá $R_L$ sobre la ganancia de la etapa? Calcula el paralelo efectivo $R_C \parallel R_L$.
    \item Calcula la ganancia de etapa cargada $A_v = v_o / v_i$.
    \item Calcula la ganancia global de tensión $G_v = v_o / v_s$.
    \item Si desconectamos $R_L$ ($R_L \to \infty$), ¿a qué valor tendería $A_v$?
    \item Si reducimos la resistencia de la fuente $R_S \to 0\,\Omega$, ¿qué ocurriría con $G_v$?
\end{enumerate}

\end{document}
